\subsection{Homological dimension of graded modules}

In this section, we assume that A is a graded ring with degrees ≥ 0. We denote by \((\mathbf{A}_{n})_{n \in \mathbf{Z}}\) its grading; we thus have \(A_{n} = 0\) for n \textless{} 0, \(A_{0}\) is a subring of A, \(J_{0} = \bigoplus_{n > 0} A_{n}\) is a two-sided ideal of A and the quotient graded ring \(A/J_{0}\) is identified with \(A_{0}\) .

Lemma 4. --- Let M be a graded A-module bounded below (X, p.~56). If \(\mathbf{A}_0\otimes_{\mathbf{A}}\mathbf{M} = 0\) , then \(\mathbf{M} = 0\) .

Since \(A_{0} \otimes_{A} M\) is isomorphic to \(M/J_{0} M\) , this is none other than II, p.~171, prop. 6.

Lemma 5. --- Let M be a graded A-module bounded below, and let

\[
s: \mathbf {A} \otimes_ {\mathbf {A} _ {0}} \mathrm{M} / \mathrm{J} _ {0} \mathrm{M} \rightarrow \mathrm{M}
\]

be a graded A-homomorphism such that \(1 \otimes_{\mathbf{A}} s: \mathbf{A}_0 \otimes_{\mathbf{A}} (\mathbf{A} \otimes_{\mathbf{A}_0} \mathbf{M} / \mathrm{J}_0 \mathbf{M}) \to \mathbf{A}_0 \otimes_{\mathbf{A}} \mathbf{M}\) is the canonical isomorphism. Then s is surjective. If \(\operatorname{Tor}_1^{\mathbf{A}}(\mathbf{A}_0, \mathbf{M}) = 0\) , s is bijective.

We have an exact sequence

\[
0 \rightarrow \operatorname{Ker} s \rightarrow A \otimes_ {A _ {0}} M / J _ {0} M \xrightarrow {s} M \rightarrow \operatorname{Coker} s \rightarrow 0
\]

and the graded A-modules Ker \(s\) and Coker \(s\) are bounded below. From the exact sequence \(\mathbf{A}_{0} \otimes_{\mathbf{A}} (\mathbf{A} \otimes_{\mathbf{A}_{0}} \mathrm{M} / \mathrm{J}_{0} \mathrm{M}) \xrightarrow{1 \otimes s} \mathbf{A}_{0} \otimes_{\mathbf{A}} \mathrm{M} \to \mathbf{A}_{0} \otimes_{\mathbf{A}} \text{Coker } s \to 0\) , we deduce that \(\mathbf{A}_{0} \otimes_{\mathbf{A}} \text{Coker } s = 0\) , hence \(s\) is surjective (lemma 4). We then have an exact sequence

\[
\operatorname{Tor} _ {1} ^ {\mathrm{A}} \left(\mathrm{A} _ {0}, \mathrm{M}\right)\rightarrow \mathrm{A} _ {0} \otimes_ {\mathrm{A}} \operatorname{Ker} s \rightarrow \mathrm{A} _ {0} \otimes_ {\mathrm{A}} \left(\mathrm{A} \otimes_ {\mathrm{A} _ {0}} \mathrm{M} / \mathrm{J} _ {0} \mathrm{M}\right) \xrightarrow {1 \otimes s} \mathrm{A} _ {0} \otimes_ {\mathrm{A}} \mathrm{M}.
\]

If \(\operatorname{Tor}_1^{\mathbf{A}}(\mathbf{A}_0, \mathbf{M}) = 0\) , then \(\mathbf{A}_0 \otimes_{\mathbf{A}} \operatorname{Ker} s = 0\) and \(s\) is injective (lemma 4).

PROPOSITION 8. --- Let M be a graded A-module bounded below.

\begin{enumerate}
\def\labelenumi{\alph{enumi})}
\tightlist
\item
  The following conditions are equivalent:
\end{enumerate}

\begin{enumerate}
\def\labelenumi{(\roman{enumi})}
\item
  M is isomorphic to a graded A-module of the form \(\mathbf{A} \otimes_{\mathbf{A}_0} \mathbf{N}\) , where \(N\) is a graded projective (resp. graded free) \(\mathbf{A}_0\) -module;
\item
  M is a projective (resp. free) graded A-module;
\item
  \(\mathbf{M} / \mathbf{J}_0\mathbf{M}\) is a projective (resp. graded free) \(\mathbf{A}_0\) -module and \(\operatorname{Tor}_1^{\mathbf{A}}(\mathbf{A}_0,\mathbf{M}) = 0\) .
\end{enumerate}

\begin{enumerate}
\def\labelenumi{\alph{enumi})}
\setcounter{enumi}{1}
\tightlist
\item
  Suppose moreover that M has a generating system consisting of homogeneous elements of bounded degrees. Then the following conditions are equivalent:
\end{enumerate}

\begin{enumerate}
\def\labelenumi{(\roman{enumi})}
\item
  The graded A-module M possesses a finite composition series whose quotients are isomorphic to graded A-modules of the form \(A \otimes_{A_{0}} N\) , where N is a flat graded \(A_{0}\) -module;
\item
  M is a flat A-module;
\item
  \(\mathbf{M} / \mathrm{J}_0\mathbf{M}\) is a flat \(\mathbf{A}_0\) -module and \(\operatorname{Tor}_1^{\mathbf{A}}(\mathbf{A}_0,\mathbf{M}) = 0\) .
\end{enumerate}

In each of the two cases, we evidently have (i) \(\Rightarrow\) (ii) \(\Rightarrow\) (iii). We must therefore prove (iii). \(\Rightarrow\) (i).

\begin{enumerate}
\def\labelenumi{\alph{enumi})}
\tightlist
\item
  The graded A-module A \(\otimes_{\mathbf{A}_0}\mathbf{M} / \mathbf{J}_0\) M is a projective A-module since \(\mathbf{M} / \mathbf{J}_0\) M is a projective \(\mathbf{A}_0\) -module. The canonical homomorphism of A-modules
\end{enumerate}

\[
p: \mathbf {M} \rightarrow \mathbf {M} / \mathrm{J} _ {0} \mathbf {M}
\]

is surjective; there therefore exists a graded A-homomorphism of degree zero

\[
s: \mathbf {A} \otimes_ {\mathbf {A} _ {0}} \mathbf {M} / \mathrm{J} _ {0} \mathbf {M} \rightarrow \mathbf {M}
\]

such that \(p \circ s(a \otimes x) = ax\) for \(a \in A\) and \(x \in M / J_0 M\) .

By Lemma 5, \(s\) is an isomorphism of A-modules from \(\mathbf{A} \otimes_{\mathbf{A}_0} \mathrm{M} / \mathrm{J}_0 \mathrm{M}\) over \(\mathbf{M}\) , whence (i).

\begin{enumerate}
\def\labelenumi{\alph{enumi})}
\setcounter{enumi}{1}
\tightlist
\item
  By the hypothesis on M, there exist integers a, b such that \(a \leqslant b\) such that M is generated by \(\oplus M_{i}\) Let us reason by induction on the positive integer b - a.
\end{enumerate}

If \(b - a = 0\) then M is generated by \(\mathbf{M}_a\) and the canonical \(\mathbf{A}_0\) -homomorphism \(\mathbf{M}_a \to \mathbf{M} / \mathbf{J}_0\) M is bijective; one then deduces from the A-homomorphism \(\mathbf{A} \otimes_{\mathbf{A}_0} \mathbf{M}_a \to \mathbf{M}\) defined by the structure of M as an A-module a graded A-homomorphism

\[
s: \mathbf {A} \otimes_ {\mathbf {A} _ {0}} \mathbf {M} / \mathrm{J} _ {0} \mathbf {M} \rightarrow \mathbf {M}
\]

satisfying the condition of Lemma 5. Then, by Lemma 5, s is bijective, whence (i).

In the general case, let \(\mathbf{M}^{(a)}\) be the (graded) sub-A-module of M generated by \(M_{a}\) and \(M'\) the quotient \(\mathbf{M}/\mathbf{M}^{(a)}\) . We have an exact sequence

\[
0 \to \mathbf {M} ^ {(a)} \xrightarrow {f} \mathbf {M} \xrightarrow {g} \mathbf {M} ^ {\prime} \to 0,
\]

whence, since \(\mathrm{Tor}_{1}^{\mathbf{A}}(\mathbf{A}_{0},\mathbf{M}) = 0\) by hypothesis, exact sequences

\begin{enumerate}
\def\labelenumi{(\arabic{enumi})}
\setcounter{enumi}{5}
\tightlist
\item
\end{enumerate}

\[
\begin{array}{c} \operatorname{Tor} _ {2} ^ {\mathsf {A}} (\mathbf {A} _ {0}, \mathbf {M} ^ {\prime}) \to \operatorname{Tor} _ {1} ^ {\mathsf {A}} (\mathbf {A} _ {0}, \mathbf {M} ^ {(a)}) \to 0 \\ 0 \to \operatorname{Tor} _ {1} ^ {\mathsf {A}} (\mathbf {A} _ {0}, \mathbf {M} ^ {\prime}) \to \operatorname{M} ^ {(a)} / \mathrm{J} _ {0} \operatorname{M} ^ {(a)} \xrightarrow {1 \otimes f} \operatorname{M} / \mathrm{J} _ {0} \operatorname{M} \xrightarrow {1 \otimes g} \operatorname{M} ^ {\prime} / \mathrm{J} _ {0} \operatorname{M} ^ {\prime} \to 0. \end{array}\tag{7}
\]

But the canonical homomorphism \(M_{a} \to M^{(a)}/J_{0} M^{(a)}\) is bijective. It follows that the homomorphism \(1 \otimes f : M^{(a)}/J_{0} M^{(a)} \to M/J_{0} M\) is injective and that its image is a sub- \(A_{0}\) -module direct summand of \(M/J_{0} M\) . It then follows from the exact sequence (7) that \(\operatorname{Tor}_{1}^{A}(A_{0}, M') = 0\) and that the \(A_{0}\) -module \(M'/J_{0} M'\) is flat since it is isomorphic to a direct summand of \(M/J_{0} M\) . By the induction hypothesis (which applies to \(M'\) since the latter is generated by the \(M'^{i}\) for \(a < i \leqslant b\) ), \(M'\) satisfies condition (i), hence is flat. One then deduces from the exact sequence (6) that \(\operatorname{Tor}_{1}^{A}(A_{0}, M^{(a)}) = 0\) , but \(M^{(a)}/J_{0} M^{(a)}\) is identified with \(M_{a}\) which is a flat \(A_{0}\) -module (as a direct summand submodule of \(M/J_{0} M\) ); by what has already been proved, the graded A-module \(M^{(a)}\) is isomorphic to \(A \otimes_{A_{0}} M_{a}\) hence also satisfies (i), which completes the proof.

COROLLARY 1. --- Let M be a finitely generated graded A-module. If the \(\mathbf{A}_0\) -module \(\mathrm{M} / \mathrm{J}_0\) M is projective (resp. graded free, resp. flat) and if \(\operatorname{Tor}_1^{\mathbf{A}}(\mathbf{A}_0, \mathbf{M}) = 0\) then the A-module M is projective (resp. graded free, resp. flat).

Corollary 2. --- Suppose that every projective \(\mathbf{A}_{0}\) -module is free (resp. that A be noetherian and that every finitely generated projective \(\mathbf{A}_{0}\) -module be free). Let M be a graded A-module bounded below (resp. a finitely generated graded A-module) and let n be an integer \(\geqslant 0\) such that \(\mathrm{dp}_{\mathbf{A}}(\mathbf{M}) \leqslant n\) . There exists an exact sequence of graded A-modules and graded homomorphisms of degree 0

\[
0 \rightarrow \mathrm{L} _ {n} \rightarrow \mathrm{L} _ {n - 1} \rightarrow \dots \rightarrow \mathrm{L} _ {0} \rightarrow \mathrm{M} \rightarrow 0
\]

where the \(L_{i}\) are graded free and bounded below (resp. graded free and finitely generated).

Indeed, there exists (X, p.~56, prop. 11) an exact sequence of graded A-modules and homomorphisms of degree 0

\[
0 \to \mathrm{L} _ {n} \to \mathrm{L} _ {n - 1} \to \dots \to \mathrm{L} _ {0} \to \mathrm{M} \to 0
\]

where the \(L_{i}\) are bounded below (resp. finitely generated) for \(0 \leqslant i \leqslant n\) and graded free for \(0 \leqslant i \leqslant n - 1\) .

Since \(\mathrm{dp}_{\mathbf{A}}(\mathbf{M}) \leqslant n\) the A-module \(L_{n}\) is projective; therefore \(\mathbf{A}_0 \otimes_{\mathbf{A}} L_n\) is a projective \(\mathbf{A}_0\) -module, hence graded free; since \(L_{n}\) is bounded below and \(\mathbf{A}_0 \otimes_{\mathbf{A}} L_n\) graded free, \(L_{n}\) is graded free (prop. 7).

COROLLARY 3 (Hilbert's syzygy theorem). --- Suppose that \(\mathbf{A}_0\) is a commutative field and that A is generated as \(\mathbf{A}_0\) -algebra by \(n\) homogeneous elements of degrees \(>0\) algebraically independent. For every graded A-module M bounded below (resp. of finite type), there exists an exact sequence of graded A-modules and graded homomorphisms of degree 0

\[
0 \rightarrow L _ {n} \rightarrow L _ {n - 1} \rightarrow \dots \rightarrow L _ {0} \rightarrow M \rightarrow 0,
\]

where the \(L_{i}\) are graded free and bounded below (resp. and finitely generated).

Indeed, \(\mathrm{dh}(\mathbf{A}) = n\) by theorem 1 of X, p.~143, and one applies cor. 2.

Remark. --- Cor. 2 also applies to the following cases:

\begin{enumerate}
\def\labelenumi{\alph{enumi})}
\item
  \(\mathbf{A}_0\) is principal and \(\mathbf{A} = \mathbf{A}_0[\mathbf{X}_1, \dots, \mathbf{X}_{n-1}]\) ;
\item
  * A \(_{0}\) is a regular local noetherian ring of dimension r and A = A \(_{0}\) {[}X \(_{1}\) , \ldots, X \(_{n}\) , r{]}. *
\end{enumerate}

COROLLARY 4. --- Assume \(\mathbf{A}_{0}\) semisimple. Let M be a graded A-module bounded below and n an integer \(\geqslant 0\) . For \(\mathrm{dp}_{\mathbf{A}}(\mathbf{M}) \leqslant n\) , it is necessary and sufficient that \(\operatorname{Tor}_{n+1}^{\mathbf{A}}(\mathbf{A}_{0}, \mathbf{M}) = 0\) .

If \(\mathrm{dp}_{\mathbf{A}}(\mathbf{M}) \leqslant n\) , then \(\operatorname{Tor}_{n+1}^{\mathbf{A}}(\mathbf{A}_0, \mathbf{M}) = 0\) (X, p.~135, lemma 1). Conversely, let \(0 \to \mathbf{K} \to \mathbf{L}_{n-1} \to \ldots \to \mathbf{L}_1 \to \mathbf{L}_0 \to \mathbf{M} \to 0\) be an exact sequence of graded A-modules bounded below such that \(\mathbf{L}_0, \ldots, \mathbf{L}_{n-1}\) are graded free (X, p.~56, prop. 11); by cor. 3 of X, p.~131, the equality \(\operatorname{Tor}_{n+1}^{\mathbf{A}}(\mathbf{A}_0, \mathbf{M}) = 0\) implies \(\operatorname{Tor}_1^{\mathbf{A}}(\mathbf{A}_0, \mathbf{K}) = 0\) ; since \(\mathbf{K} / \mathbf{J}_0\mathbf{K}\) is a projective \(\mathbf{A}_0\) -module since \(\mathbf{A}_0\) is semisimple (X, p.~140, prop. 6), K is projective by prop. 8 (X, p.~144), and \(\mathrm{dp}_{\mathbf{A}}(\mathbf{M}) \leqslant n\) .

COROLLARY 5. --- Suppose the ring \(\mathbf{A}_0\) semisimple. If \(\operatorname{Tor}_{n+1}^{\mathbf{A}}(\mathbf{A}_0, \mathbf{A}_0) = 0\) , we have \(\mathrm{dp}_{\mathbf{A}}(\mathbf{M}) \leqslant n\) for every graded \(\mathbf{A}\) -module bounded below.

Let A° denote the opposite graded ring of A: we have (A°)₀ = (A₀)°, so (A°)₀ is semisimple (VIII, § 5, n° 1, remark 3). Since \(\operatorname{Tor}_{n+1}^{A°}(A₀°, A₀°) = 0\) , we have \(\mathrm{dp}_A(A₀s) \leqslant n\) by Cor. 4: this implies \(\operatorname{Tor}_{n+1}^{A°}(M₀, A₀°) = 0\) for every A-module M, hence \(\operatorname{Tor}_{n+1}^{A}(A₀, M) = 0\) and we apply Cor. 4.

